\documentclass[11pt,a4paper]{article}
\usepackage[T1]{fontenc}
\usepackage{lmodern,amsmath,amssymb,amsthm,mathtools,booktabs,microtype}
\usepackage[margin=27mm]{geometry}
\usepackage[colorlinks=true,linkcolor=blue,citecolor=blue,urlcolor=blue]{hyperref}
\setlength{\emergencystretch}{2em}
\newtheorem{theorem}{Theorem}[section]
\newtheorem{lemma}[theorem]{Lemma}
\newtheorem{proposition}[theorem]{Proposition}
\newtheorem{corollary}[theorem]{Corollary}
\newcommand{\Q}{\mathbb Q}\newcommand{\Z}{\mathbb Z}
\DeclareMathOperator{\Tr}{Tr}\DeclareMathOperator{\N}{N}
\title{Bounded complement levels from native detector norms}
\author{}\date{16 November 2026}
\begin{document}\maketitle
\begin{abstract}
We construct complement levels as fixed rational functions of the native
elliptic detector. At every odd valuation of its physical weight, the
factor $\N(W)/(1+\N(W))$ is a local unit. The level
$17(\N(W)/(1+\N(W)))^2$ therefore preserves the negative-odd-valuation
implication of the constant-level complement test. Its denominator
and complement are units at every character of a prepared actual point
algebra. Consequently the same finite-algebra encoding gives a
twenty-one-variable or nine-variable conditional reduction, with no
new retained parameter. The missing premise is simultaneous rational
realization with the actual prescribed first coordinates
$2W_0/(1+\N(W))$. These coordinates lie on an explicit rational quadric.
No actual uniform realization or new variable-count record is asserted.
\end{abstract}

\section{Native products and parity}

Let $K=\Q(\rho)$ with $\rho^2=2$. A bar denotes conjugation.
Let $S$ be the fixed exceptional prime set of the chosen source;
in particular $2,3,17\in S$. Use the general source or either
conditional source covered by~\cite{Complements}, retaining its
arithmetic premises and every point-stage guard.
The general source with the norm-compatible scalar of~\cite{NormPreparation}
is included with its specified fixed exceptional set.

In each block the native detector has the form
\begin{equation}\label{eq:native}
 W=\prod_{k=1}^r(1+w h_k^2),\qquad
 N=\N_{K/\Q}(W),\quad V=\Tr_{K/\Q}(W^{-1}),\qquad
 b=\frac{4w}{V^2},
\end{equation}
where $w\in\Q^\times$, the $h_k\in K$ are rational expressions in
the retained point data, and $NV\ne0$.
For the general source, $w=\rho_*R$ and the factors are exactly
\[
 1+w\left(\frac{C_i}{\delta w^{a_i}}\right)^2,
 \qquad
 1+w\left(\frac{G_k}{w^{\kappa_k}}\right)^2.
\]
Here $\rho_*$ denotes the scalar coefficient, not the basis element
$\rho=\sqrt2$.
For a conditional source,
\[
 w=\frac{F}{4\delta^8},\qquad
 1+\frac{\delta^4C_i^2}{F}
 =1+w\left(\frac{C_i}{2\delta^2w}\right)^2.
\]
Thus~\eqref{eq:native} is an exact presentation of the unchanged
native product, not a replacement of its coefficient input or trace.

\begin{lemma}\label{lem:parity}
Fix $p\notin S$ with $v_p(w)$ odd, of either sign.
At every place $\mathfrak p\mid p$, normalized by
$v_{\mathfrak p}(p)=1$, one has $v_{\mathfrak p}(W)\le0$.
Consequently $v_p(N)\le0$. If $v_p(N)=0$, then
\begin{equation}\label{eq:N-one}
                              N\equiv1\pmod p.
\end{equation}
\end{lemma}
\begin{proof}
The extension is unramified at $p$, so local valuations of elements
of $K$ are integers. A nonzero second summand $wh_k^2$ has odd
valuation. Its valuation cannot equal that of $1$, and hence
\[
 v_{\mathfrak p}(1+wh_k^2)
       =\min\{0,v_p(w)+2v_{\mathfrak p}(h_k)\}\le0.
\]
A zero $h_k$ contributes the factor $1$.
The norm valuation is the sum of these nonpositive factor valuations,
weighted by residue degrees. If it is zero, every summand is zero.
Each nonzero $wh_k^2$ then has positive valuation, so every factor
reduces to $1$ and $W\equiv1$ at every place above $p$.
Taking the norm gives~\eqref{eq:N-one}, at split and inert primes alike.
\end{proof}

\begin{corollary}\label{cor:residues}
At the primes in Lemma~\ref{lem:parity}, $W^{-1}$ is locally integral
and reduces to either $0$ or $1$ at each place. The pair
\[
            \bigl(\Tr(W^{-1}),\N(W^{-1})\bigr)
\]
therefore reduces to one of $(0,0),(1,0),(2,1)$.
At inert primes only $(0,0)$ and $(2,1)$ occur.
\end{corollary}
\begin{proof}
If some factor of $W$ has negative valuation, its inverse and hence
$W^{-1}$ reduce to zero. Otherwise all factors reduce to one.
At a split prime the trace and norm are the sum and product of the
two residues. At an inert prime the two conjugate residues agree,
because $0$ and $1$ belong to the prime field.
\end{proof}

\section{A level with the required local implication}

Fix $c\in\Z[S^{-1}]$ with $c\ge0$ and
\begin{equation}\label{eq:c}
                        1+c\in\Z[S^{-1}]^\times.
\end{equation}
The choices $c=0$ and $c=1$ are always allowed. Let $M$ be a fixed
positive integer and put
\begin{equation}\label{eq:level}
              m=\frac{MN}{c+N},\qquad a=17m^2,
\end{equation}
with the added denominator guard $c+N\ne0$.
These are rational functions of the native point data; neither $m$
nor $a$ is a new retained coordinate.

\begin{lemma}\label{lem:unit}
If $p\notin S$ and $v_p(w)$ is odd, then
\[
                   v_p\left(\frac{N}{c+N}\right)=0,
                   \qquad v_p(a)=2v_p(M)\ge0.
\]
\end{lemma}
\begin{proof}
If $v_p(N)<0$, then $v_p(c)\ge0$ and $N$ determines the valuation
of $c+N$. If $v_p(N)=0$, Lemma~\ref{lem:parity} and
\eqref{eq:c} give $c+N\equiv c+1\not\equiv0\pmod p$.
These are the only cases. Also $17$ is a unit at $p$.
\end{proof}

\begin{proposition}\label{prop:local}
Let $d\in\Q^\times$ have even valuations outside $S$.
If the actual native values satisfy
\begin{equation}\label{eq:complement}
                                q^2=d\,a(a-b)
\end{equation}
for some $q\in\Q$, then $b$ has no negative odd valuation
outside $S$.
\end{proposition}
\begin{proof}
Since $b=4w/V^2$, an odd valuation of $b$ is an odd valuation
of $w$. If $v_p(b)<0$ is odd, Lemma~\ref{lem:unit} gives
$v_p(a)\ge0$ even, so $v_p(a-b)=v_p(b)$.
The right side of~\eqref{eq:complement} then has odd valuation,
contradicting squareness. If $a=b$, the same comparison already
contradicts a negative odd valuation of $b$.
\end{proof}

The level $a$ need not be admissible at all primes outside $S$.
Its nonnegative even valuation is required only at primes where the
native weight has odd valuation, and Lemma~\ref{lem:unit} supplies
exactly that condition. In particular the proof includes negative
odd valuations of the general scalar and uses no assumption that
a rational input is already integral.

\section{All characters and the finite encoding}

Let $\mathcal B$ be the tensor product of the native point algebras
at an actual guarded point-stage specialization. By
\cite[Lemma~1.1]{Complements}, it is split by $K$ and every geometric
character has image in this totally real field.
The norm and trace in~\eqref{eq:native} are first computed in
$\mathcal B\otimes_\Q K$, with conjugation acting on the $K$ factor.
Thus $N,V$ and $b$ belong to the guarded localization of $\mathcal B$.

\begin{lemma}\label{lem:guards}
Suppose the native inputs are prepared as in~\cite{Complements}.
Then $c+N$, $m$, and $a-b$ are units of the whole point algebra.
In particular every additional denominator and complement norm
required by~\eqref{eq:level}--\eqref{eq:complement} is nonzero.
\end{lemma}
\begin{proof}
At the selected rational character the prepared input is positive,
so $w>0$ by~\eqref{eq:native}; this rational scalar is independent
of the character. Extend any geometric character of $\mathcal B$
to the two real embeddings of its tensor product with $K$.
Every factor $1+wh_k^2$ is at least one in each such real embedding.
Hence the image of $N$ is at least one, and $c+N>0$.
The native guards already make $N$ a unit, so $m$ is a unit too.

Write $b=wT^2$ in the point algebra, where $T=2/V$ is a unit.
Preparation at the selected rational character implies that
$v_{17}(w)$ is even. If a geometric character $\chi$ annihilated
$a-b$, it would give
\[
                       \frac{17}{w}
                         =\left(\frac{\chi(T)}{\chi(m)}\right)^2
                         \in K^{\times2}.
\]
A nonzero rational square in $K$ has rational square class $1$
or $2$. Such a number has even valuation at $17$, whereas
$v_{17}(17/w)$ is odd. This contradiction proves that $a-b$
is a unit. The argument concerns every geometric character,
not just the selected rational one.
\end{proof}

Fix finite nonempty sets of coefficients $M$, $c$, and $d$ as above,
with $d>0$ and its square class supported on $S$.
They are fixed data of the polynomial, not quantified parameters.

\begin{theorem}\label{thm:reduction}
Fix a coefficient member of the general source or either conditional
source in its sufficient soundness range, retaining its arithmetic
premises. Suppose that for every integer input in its completeness
range there is an actual prepared point-stage tuple satisfying all
inherited guards and, in each of the two blocks, equation
\eqref{eq:complement} for one of the fixed coefficient choices.
Then its pole stage can be replaced without adding any retained
coordinate, and the single-coordinate finite-algebra encoding gives
twenty-one total variables for the general source or nine for either
conditional source. The input is included.
The stated simultaneous actual realization is an additional
unproved premise.
\end{theorem}
\begin{proof}
For each fixed pair of coefficient choices, form $a_i$ and $b_i$
as rational functions in the existing free rank-$4^p$ point algebra,
where $p=9$ or $p=3$. The new expressions use only rational
operations on that algebra. Clear their denominators, adding the
norms of the numerators of $c_i+N_i$ to the original scalar guard.
Represent
\[
                        a_i(a_i-b_i)=R_i/D_i
\]
in the integral point algebra with $D_i$ a guarded unit.
Choose each square-class representative $d_i$ to be a positive
squarefree integer and adjoin the two monic roots
\[
                            Q_i^2=d_iR_iD_i.
\]
At a rational character, $Q_i/D_i$ is precisely the root in
\eqref{eq:complement}. Include the complement norms among the
scalar guards. Lemma~\ref{lem:guards} makes all these additions
nonzero on the actual prepared tuples supplied by the hypothesis.
The two quadratic radicands are units and the enlarged algebra
has rank $4^{p+1}$. The integral denominator envelopes, multiplication
matrices and one-coordinate separator of~\cite{Unit} therefore apply.
Neither the norm $N_i$, the rational function $m_i$, its denominator,
nor the finite roots introduce an unrestricted rational coordinate.

For soundness, Proposition~\ref{prop:local} supplies the same local
negative-odd-valuation implication used by the native source.
Its remaining point, local extraction and height arguments are
unchanged. Completeness follows from the stated actual realization
and Lemma~\ref{lem:guards}, followed by the finite-algebra separator.
The retained coordinates are the input, $2p$ point coordinates and
the slope. One encoder gives $2p+3$, namely $21$ or $9$.
As in the parent conditional source, a single input factor includes
zero without another coordinate.

Finally take a product of these defining polynomials over the fixed
finite coefficient choices, using the same variable names. Its
rational zero set is the union of the separately guarded alternatives.
This does not increase the variable count.
\end{proof}

The choice $c=0$ recovers the constant levels $17M^2$.
Theorem~\ref{thm:reduction} therefore permits fixed finite families
containing both those levels and the nonconstant native levels.
No degree or coefficient bound for the resulting encodings is
assigned, and no record count follows without the realization premise.

\section{The prescribed-coordinate norm equations}

Write $W=W_0+W_1\rho$ at the selected rational character.
All quantities in the following equations are the actual native ones.
From~\eqref{eq:native} and~\eqref{eq:level},
\[
                           mV=\frac{2MW_0}{c+N}.
\]
Because $mV\ne0$, equation~\eqref{eq:complement} is equivalent to
\begin{equation}\label{eq:normtarget}
                  \left(\frac{2MW_0}{c+N}\right)^2-dy^2
                             =\frac{4w}{17}.
\end{equation}
Indeed the mutually inverse witness substitutions are
\[
                          y=\frac{qV}{17dm},\qquad
                          q=\frac{17dmy}{V}.
\]
Thus the new rational-level theorem prescribes a different first
coordinate on the same bare norm conic; it does not select that
coordinate freely.

For the canonical choice $M=c=1$, put
\begin{equation}\label{eq:cayley}
 u=\frac{2W_0}{1+N},\qquad
 v=\frac{2W_1}{1+N},\qquad
 z=\frac{1-N}{1+N}.
\end{equation}
Direct expansion gives
\begin{equation}\label{eq:quadric}
                              u^2-2v^2+z^2=1.
\end{equation}
The actual norm-incidence target is now
\[
                     u_i^2-d_i y_i^2=4w_i/17\qquad(i=0,1),
\]
with $(u_i,v_i,z_i)$ fixed by~\eqref{eq:cayley} and by the coupled
native arguments. Bare norm-conic solubility does not prescribe these
coordinates.

\begin{proposition}\label{prop:bounds}
At every prepared actual character,
\[
             \frac{17}{4}\le a<17,\qquad 0<u\le1,
             \qquad -1<z\le0
\]
for the canonical choice. More generally the levels in
\eqref{eq:level} lie between $17M^2/(1+c)^2$ and $17M^2$.
\end{proposition}
\begin{proof}
The two real values $r=W_0+\rho W_1$ and $s=W_0-\rho W_1$
are at least one. Thus $N=rs\ge1$ and
\[
                 1+rs-r-s=(r-1)(s-1)\ge0.
\]
These inequalities give $0<u=(r+s)/(1+rs)\le1$,
$-1<z=(1-rs)/(1+rs)\le0$, and
$1/2\le N/(1+N)<1$.
For general $c\ge0$, one has
$1/(1+c)\le N/(c+N)\le1$; squaring proves the final assertion.
\end{proof}

For a conditional source, increasing the prescribed real arguments
on the fixed compact point neighborhoods gives $W_i\to1$ and
$b_i\to0$. Its canonical levels therefore tend to $17/4$.
For the general source, the parent preparation first makes the scalar
factors as close to one as required, before fixing the root and scale
data. Increasing the real slope then makes the comparison factors and
$R_i$ as close to one as required. Thus $b_i$ can be made arbitrarily
close to the fixed $\rho_*$, where $0<\rho_*<1$, while the levels
are as close to $17/4$ as required. These ordered choices give
$b_i<a_i$ in both sources. They check the real inequality for positive
complement classes; they supply no simultaneous rational square roots
in~\eqref{eq:normtarget}.

\section{Exact identities and scope}

The certificate in
\nolinkurl{computation/pole_parameters/bounded_norm_levels_v48}
verifies the conditional native factor presentation, both directions
of the complement/norm substitution, the quadric identity by two
polynomial calculations, the real-bound factorization, and the
possible split and inert inverse-trace residues. A supplementary
valuation screen includes positive and negative odd weight orders.
The all-valuation and whole-character assertions are proved above.
The certificate does not assert actual uniform rational realization.

\begin{thebibliography}{9}
\bibitem{Complements}\emph{Unit guards and a finite square-complement reduction for elliptic detectors},
\nolinkurl{papers/pole_parameters/guarded_complement_reduction_20261116_v45.tex}.
\bibitem{FiniteClasses}\emph{Finite square-class conditions for native detector inputs},
\nolinkurl{papers/pole_parameters/finite_complement_classes_20261116_v46.tex}.
\bibitem{Unit}\emph{A single norm determinant for the 25- and 13-variable constructions},
\nolinkurl{papers/finite_algebra/unit_norm_20261114_v28.tex}.
\bibitem{NormPreparation}\emph{Norm preparation for general and conditional elliptic detectors},
\nolinkurl{papers/pole_parameters/norm_compatible_scalar_20261116_v47.tex}.
\end{thebibliography}
\end{document}
